% !TEX root = ../thesis.tex
% 人工智能生成合成内容标识 — 哈工大写作指南 §2.7
% 对应《关于印发<人工智能生成合成内容标识办法>的通知》(国家网信办, 2025)

\chapter*{人工智能生成合成内容标识}
\addcontentsline{toc}{chapter}{人工智能生成合成内容标识}
\markboth{人工智能生成合成内容标识}{人工智能生成合成内容标识}

依照《关于印发\textless 人工智能生成合成内容标识办法\textgreater 的通知》及哈尔滨工业大学《研究生学位论文写作指南（理工类）》（2025 年 3 月）\S 2.7 的要求，
本论文在撰写过程中对人工智能（AI）工具的使用情况作如下说明：

\section*{使用范围}

\begin{itemize}
  \item \textbf{文献调研辅助}：在课题立项与文献整理阶段，
        作者使用了大语言模型（OpenAI ChatGPT、Anthropic Claude 等）辅助进行文献关键词扩展、
        概念解释和英文表达润色，所有参考文献的最终确认与引用均由作者完成。

  \item \textbf{代码工程辅助}：本论文涉及的代理模型训练脚本
        （\texttt{optimization\_v2/tools/sweep\_moe.py} 等）、
        数据转换流水线和可视化代码片段，
        部分采用大语言模型作为编程助手协助完成模板生成与代码补全；
        所有代码的最终功能正确性、数值结果与物理含义均经作者独立验证。

  \item \textbf{论文写作辅助}：在论文撰写阶段，
        AI 工具用于段落组织、英文摘要润色、表格 LaTeX 模板生成等纯形式性的辅助工作；
        所有研究内容、数据、图表、公式推导、实验结论及创新观点均由作者本人提出与撰写。
\end{itemize}

\section*{未使用范围}

下列内容\textbf{未使用任何 AI 生成合成技术}：
\begin{itemize}
  \item 所有数值实验数据（包括 LLFVW 仿真数据、代理模型训练 / 测试结果、CFD 验证数据）；
  \item 所有论文插图（除工程示意流程图外，均由作者基于实验数据绘制）；
  \item 课题的研究方法、技术路线与创新结论；
  \item 引用的参考文献内容。
\end{itemize}

\section*{标识承诺}

作者承诺，对本论文中涉及 AI 工具辅助的部分已尽到合理审查与验证义务，
所有 AI 生成合成内容均已用于规范的辅助用途，
不存在以 AI 生成内容替代独立研究工作的情形。
对本论文最终内容的真实性、原创性与学术规范性，
作者承担全部责任。

\vspace{2em}
\begin{flushright}
作者签名：\underline{\hspace{4cm}}\quad\quad 日期：\underline{\hspace{3cm}} \\[1em]
导师签名：\underline{\hspace{4cm}}\quad\quad 日期：\underline{\hspace{3cm}}
\end{flushright}
